\section{时间推进实现}

本章说明一个时间步内代码具体做了什么。对应的主函数是：
\begin{lstlisting}
advance_one_maccormack_step(&solver);
\end{lstlisting}
它里面的主要顺序如下：
\begin{lstlisting}
solver->time.dt = compute_dt_from_cfl(solver);
apply_artificial_viscosity_filter(solver);
compute_flux_from_conserved(solver, &solver->state.filtered,
                            &solver->flux.filtered);
compute_maccormack_predictor(solver, stencil);
compute_flux_from_conserved(solver, &solver->state.predicted,
                            &solver->flux.predicted);
compute_maccormack_corrector(solver, stencil);
\end{lstlisting}

\subsection{CFL 时间步}

显式格式的时间步按最大波速决定：
\[
\Delta t = CFL\frac{\Delta x}{\max_i(|u_i|+a_i)}.
\]
代码先从当前守恒量反算 \(\rho,u,p,E,a\)，然后扫描所有网格点求
\(\max_i(|u_i|+a_i)\)。如果最后一步会超过给定的终止时间，
就把 \(\Delta t\) 截短到刚好到达 \(t_{\max}\)。

\subsection{人工粘性滤波}

根据课堂给出的滤波公式，先对当前守恒量做一次平滑：
\[
\bar{\bm Q}^{\,n}_i
=\bm Q_i^n
+\frac12\nu_i
\left(\bm Q_{i+1}^n-2\bm Q_i^n+\bm Q_{i-1}^n\right).
\]
其中
\[
\nu_i=CFL_{\mathrm{actual}}(1-CFL_{\mathrm{actual}})\beta\theta_i.
\]

因为最后一个时间步可能被截短，所以这里用的是实际 CFL：
\[
CFL_{\mathrm{actual}}
=\max_i(|u_i|+a_i)\frac{\Delta t}{\Delta x}.
\]

老师要求保留三种人工粘性算法，所以 \(\theta_i\) 可以用
\(\rho\)、\(u\) 或 \(p\) 来算。统一写成 \(s_i\)：
\[
\theta_i=
\frac{|s_{i+1}-2s_i+s_{i-1}|}
{|s_{i+1}-s_i|+|s_i-s_{i-1}|+\epsilon}.
\]
代码中通过输入选项选择 sensor：
\begin{lstlisting}
VISCOSITY_SENSOR_RHO = 1
VISCOSITY_SENSOR_U   = 2
VISCOSITY_SENSOR_P   = 3
\end{lstlisting}

如果关闭人工粘性，程序直接令
\(\bar{\bm Q}^{\,n}=\bm Q^n\)，后面的 MacCormack 步骤仍然可以正常执行。

\subsection{两种 MacCormack 写法}

本程序支持两种常见写法。第一种是预测步用 FTBS，校正步用 FTFS：
\[
\tilde{\bm Q}^{\,n+1}_i
=\bar{\bm Q}^{\,n}_i
-\lambda\left[
\bm F(\bar{\bm Q}^{\,n}_i)
-\bm F(\bar{\bm Q}^{\,n}_{i-1})
\right],
\]
\[
\bm Q^{n+1}_i
=\frac12\left(
\bar{\bm Q}^{\,n}_i+\tilde{\bm Q}^{\,n+1}_i
\right)
-\frac{\lambda}{2}
\left[
\bm F(\tilde{\bm Q}^{\,n+1}_{i+1})
-\bm F(\tilde{\bm Q}^{\,n+1}_i)
\right].
\]

第二种是预测步用 FTFS，校正步用 FTBS：
\[
\tilde{\bm Q}^{\,n+1}_i
=\bar{\bm Q}^{\,n}_i
-\lambda\left[
\bm F(\bar{\bm Q}^{\,n}_{i+1})
-\bm F(\bar{\bm Q}^{\,n}_i)
\right],
\]
\[
\bm Q^{n+1}_i
=\frac12\left(
\bar{\bm Q}^{\,n}_i+\tilde{\bm Q}^{\,n+1}_i
\right)
-\frac{\lambda}{2}
\left[
\bm F(\tilde{\bm Q}^{\,n+1}_i)
-\bm F(\tilde{\bm Q}^{\,n+1}_{i-1})
\right].
\]

其中
\[
\lambda=\frac{\Delta t}{\Delta x}.
\]

\subsection{固定使用和交替使用}

程序运行时可以选择四种模式：
\begin{enumerate}[label=\arabic*.]
    \item 每一步都使用 FTBS/FTFS；
    \item 每一步都使用 FTFS/FTBS；
    \item 两种写法交替使用，并且第一步用 FTBS/FTFS；
    \item 两种写法交替使用，并且第一步用 FTFS/FTBS。
\end{enumerate}

交替模式通过当前步数的奇偶性判断本步使用哪一种写法。
这样可以比较单一写法和交替写法对数值结果的影响。

\subsection{边界条件}

当前边界条件为简单的零梯度外推：
\[
\bm Q_0=\bm Q_1,\qquad
\bm Q_{N_x-1}=\bm Q_{N_x-2}.
\]
滤波状态、预测状态和校正状态都会各自处理边界，避免差分时读到没有定义的边界值。
